\uuid{E9wX}
\titre{Ensemble}
\auteur{bodin ; hanani ; mzari}
\organisation{exo7}
\licence{CC BY-NC-SA 4.0}
\datecreate{2020-01-20}
\niveau{L1}
\module{Algèbre}
\chapitre{Logique, ensemble, raisonnement}
\sousChapitre{Ensemble}
\difficulte{1}
\theme{Ensembles, applications, ensemble}
\exo7qcmid{lille-1.ensembles.001}

\contenu{
\begin{question}
\qclassification{100.02,101.01,102.01,102.02}
\qtags{motcle=ensemble}
Soit $A=\{x\in \Rr\mid (x+8)^2=9^2\}$. Sous quelle forme peut-on encore écrire l'ensemble $A$ ?
\begin{answers}
    \bad{$A=\{1\}$}
    \bad{$A=\varnothing$}
    \bad{$A=\{-17\}$}
    \good{$A=\{1,-17\}$}
\end{answers}
\begin{explanations}
Les éléments de $A$ sont les solutions de l'équation $(x+8)^2=9^2$, c'est-à-dire $1$ et $-17$.
\end{explanations}
\end{question}
}
